\documentclass[acmsmall,screen,review,anonymous]{acmart}
\settopmatter{printfolios=true,printccs=false,printacmref=false}

\setcopyright{none}

%% NOTE that a single column version is required for
%% submission and peer review. This can be done by changing
%% the \doucmentclass[...]{acmart} in this template to
%% \documentclass[manuscript,screen]{acmart}
%%
%% To ensure 100% compatibility, please check the white list of
%% approved LaTeX packages to be used with the Master Article Template at
%% https://www.acm.org/publications/taps/whitelist-of-latex-packages
%% before creating your document. The white list page provides
%% information on how to submit additional LaTeX packages for
%% review and adoption.
%% Fonts used in the template cannot be substituted; margin
%% adjustments are not allowed.
%%
%% \BibTeX command to typeset BibTeX logo in the docs
\AtBeginDocument{%
  \providecommand\BibTeX{{%
    \normalfont B\kern-0.5em{\scshape i\kern-0.25em b}\kern-0.8em\TeX}}}

%% Rights management information.  This information is sent to you
%% when you complete the rights form.  These commands have SAMPLE
%% values in them; it is your responsibility as an author to replace
%% the commands and values with those provided to you when you
%% complete the rights form.
% \setcopyright{acmcopyright}
% \copyrightyear{2018}
% \acmYear{2018}
% \acmDOI{XXXXXXX.XXXXXXX}


%%
%% These commands are for a JOURNAL article.
% \acmJournal{JACM}
% \acmVolume{37}
% \acmNumber{4}
% \acmArticle{111}
% \acmMonth{8}

%%
%% Submission ID.
%% Use this when submitting an article to a sponsored event. You'll
%% receive a unique submission ID from the organizers
%% of the event, and this ID should be used as the parameter to this command.
%%\acmSubmissionID{123-A56-BU3}

%%
%% For managing citations, it is recommended to use bibliography
%% files in BibTeX format.
%%
%% You can then either use BibTeX with the ACM-Reference-Format style,
%% or BibLaTeX with the acmnumeric or acmauthoryear sytles, that include
%% support for advanced citation of software artefact from the
%% biblatex-software package, also separately available on CTAN.
%%
%% Look at the sample-*-biblatex.tex files for templates showcasing
%% the biblatex styles.
%%

%%
%% The majority of ACM publications use numbered citations and
%% references.  The command \citestyle{authoryear} switches to the
%% "author year" style.
%%
%% If you are preparing content for an event
%% sponsored by ACM SIGGRAPH, you must use the "author year" style of
%% citations and references.
%% Uncommenting
%% the next command will enable that style.
%\citestyle{acmauthoryear}

%%
%% end of the preamble, start of the body of the document source.

\usepackage{hyperref, wrapfig, amsthm, pict2e, booktabs, subcaption,
  subfiles, rotating, pifont, indentfirst, amsmath, amsfonts,
  mathrsfs, adjustbox, mathtools,booktabs, microtype, ebproof,
  enumitem, multirow, multicol, setspace, xcolor, lineno, listings,
  stmaryrd, xspace, stackengine, makecell}

\usepackage{empheq}
\newcommand*\widefbox[1]{\fbox{\hspace{2em}#1\hspace{2em}}}
\usepackage[tableposition=top]{caption}
\usepackage{tikz-cd}
\usetikzlibrary{decorations.pathmorphing}
\usepackage[ruled,linesnumbered,vlined]{algorithm2e}
\newcommand\mycommfont[1]{\footnotesize\ttfamily\textcolor{DeepGreen}{#1}}
\SetCommentSty{mycommfont}
\usepackage{graphicx}
\urlstyle{same}
\usepackage{minted}
\usemintedstyle{borland}
\usepackage{tikz}
\newcounter{markeq}
\setcounter{markeq}{0}
\usepackage{spacingtricks}
\newcommand*\tstrut[1]{\vstrut{#1}}
\newcommand*\bstrut[1]{\vstrut[#1]{0pt}}
\usepackage[T1]{fontenc}
\newcommand*\mystrut[1]{\vrule width0pt height0pt depth#1\relax}
% \usepackage{bm}

\newcommand{\tyName}{\mbox{\textsf{u}\textsc{Hat}}\xspace}

% ZZ's Macros
\usepackage{spacingtricks, commands, refinementtydef}

\begin{document}

%%
%% The "title" command has an optional parameter,
%% allowing the author to define a "short title" to be used in page headers.
% \title{Synthesizing Test Controllers from Types: Property-Guided Bug-Finding for Distributed System Models}
% Optional [ ] = running header title: keeps ``Strategy Types'' visible on every page.
%\title[Context Types]{Context Types: Unifying May- and Must-Style Reasoning through Resource-Modal Typing}

\title[Context Types]{Context Types: Bunched Modalities for Unifying Safety and Reachability}
%%
%% The "author" command and its associated commands are used to define
%% the authors and their affiliations.
%% Of note is the shared affiliation of the first two authors, and the
%% "authornote" and "authornotemark" commands
%% used to denote shared contribution to the research.
\author{Zhe Zhou}
%\authornote{with author1 note}          %% \authornote is optional;
                                        %% can be repeated if necessary
%\orcid{nnnn-nnnn-nnnn-nnnn}             %% \orcid is optional
\affiliation{
  \institution{Purdue University}            %% \institution is required
  \country{USA}                    %% \country is recommended
}
% \email{first1.last1@inst1.edu}          %% \email is recommended

\author{Benjamin Delaware}
\affiliation{
  \institution{Purdue University}            %% \institution is required
  \country{USA}                    %% \country is recommended
}

\author{Suresh Jagannathan}
\affiliation{
  \institution{Purdue University}            %% \institution is required
  \country{USA}                    %% \country is recommended
}

%%
%% By default, the full list of authors will be used in the page
%% headers. Often, this list is too long, and will overlap
%% other information printed in the page headers. This command allows
%% the author to define a more concise list
%% of authors' names for this purpose.
% \renewcommand{\shortauthors}{Trovato and Tobin, et al.}

%%
%% The abstract is a short summary of the work to be presented in the
%% article.


\begin{abstract}
  Most type systems are designed to verify safety properties: the
  judgement $e : \tau$ holds if $e$ \emph{may} evaluate to any value
  contained in the denotation of $\tau$, but nothing more. More recent
  type systems have considered the complementary class of reachability
  properties: in these systems, the typing judgment holds if $e$
  \emph{must} be able to evaluate to every value denoted by $\tau$,
  and possibly more. This paper introduces a refinement type system
  for higher-order functional programs where types can capture both
  safety and reachability properties. Simultaneously supporting both
  kinds of specification is challenging because each confers different
  reasoning capabilities when they appear in a typing context. We
  address this challenge by treating typing contexts as a kind of
  resource that mediates which capabilities are available when typing
  a program. We formalize these notions by developing a resource
  algebra of contextual capabilities and a bunched modal context logic
  over this algebra. We then interpret our type system in this logic,
  establishing soundness with respect to a denotational semantics that
  unifies safety and reachability guarantees within a single typing
  judgement.
  % We further extend our type system to support nondeterminism by
  % introducing a persistency modality that governs resource
  % duplication.
\end{abstract}

\maketitle

\subfile{sections/1-intro}
\subfile{sections/2-overview}
\subfile{sections/3-typing}
\subfile{sections/4-logic}
\subfile{sections/5-extension}
\subfile{sections/6-case}
\subfile{sections/7-discussion}
\subfile{sections/8-related}
\subfile{sections/9-conclusion}
% \subfile{sections/3-rules}

\bibliographystyle{ACM-Reference-Format}
\bibliography{bibliography}

% \section*{Acknowledgements}

%%
%% The next two lines define the bibliography style to be used, and
%% the bibliography file.

% \bibliographystyle{ACM-Reference-Format}
% \bibliographystyle{plainnat}
% \bibliographystyle{plain}
% \bibliographystyle{ACM-Reference-Format}
% \citestyle{acmauthoryear}
% \bibliography{bibliography}

\newpage
\appendix
\subfile{sections/tech/outlines}
\subfile{sections/tech/1-semantics}
\subfile{sections/tech/2-basictyping}
\subfile{sections/tech/3-pre}

\end{document}
\endinput
%%
%% End of file `sample-acmsmall.tex'.
